\section{Related Work}
\label{sec:related}

\subsection{Linear Forecasters}
\label{sec:related:linear}

The LTSF-Linear family \citep{zeng2023dlinear} showed that embarrassingly simple
linear models match or beat elaborate Transformer forecasters on the standard
long-term benchmarks. Three variants are direct baselines for us. \emph{Linear}
is a single linear layer mapping a length-$L$ lookback to a length-$H$ horizon,
applied per channel. \emph{NLinear} subtracts the last value of the lookback
before the linear map and adds it back afterwards---a one-line normalization
that handles simple distribution shift between train and test. \emph{DLinear}
decomposes the series into a moving-average trend and a seasonal/remainder
component (the additive decomposition introduced by Autoformer
\citep{wu2021autoformer}), forecasts each with its own linear layer, and sums
them. RLinear \citep{li2023rlinear} analyzes \emph{why} linear maps work and
introduces RLinear, the combination of RevIN \citep{kim2022revin} with a single
linear layer; its central findings are that the linear map captures periodicity
and that RevIN and channel independence (CI) are the two ingredients that make
linear models robust.

DLinear and RLinear are \method's direct competitors and primary baselines. The
key conceptual contrast is the \emph{domain of decomposition}. DLinear
decomposes in the time domain via a fixed moving-average kernel into trend and
remainder---a fixed low-pass smoothing whose remainder is everything the average
misses, forecast by an unconstrained head. \method\ decomposes in the frequency
domain via a \emph{learnable} spectral filter into low- and high-frequency
bands; a moving average is one particular fixed low-pass filter, and \method\
generalizes it by learning the band split end-to-end and giving each band its
own lightweight head. Against RLinear, \method\ shares the RevIN+linear backbone
but (a) replaces the single linear map with a frequency-decomposed set of heads
and (b) replaces vanilla RevIN with the horizon-adaptive \arevin\
(Section~\ref{sec:method:arevin}). The ablation that drops the frequency split
and closes the \arevin\ gate recovers an RLinear-like model, the controlled
comparison that isolates our contributions
(Section~\ref{sec:ablation}).

\subsection{Non-stationarity and Reversible Normalization}
\label{sec:related:norm}

Real series are non-stationary: mean and variance drift between the lookback and
the horizon, and between train and test. RevIN \citep{kim2022revin} is a
symmetric, instance-wise normalize-then-denormalize wrapper with a learnable
affine transform: it removes each instance's mean and standard deviation at the
input and restores them at the output, and it is now a near-universal plug-in.
Non-stationary Transformers \citep{liu2022nonstationary} make a subtler point:
naive stationarization can \emph{over-stationarize} and erase the very signal a
model needs, so they pair series stationarization with a de-stationary attention
that re-injects the removed statistics into the attention computation.

This is \method's clearest novelty. Vanilla RevIN applies a single,
horizon-agnostic de-normalization, adding back exactly the lookback mean and
variance to \emph{every} forecast step; but under distribution shift the
appropriate level and scale typically drift across the horizon, so the lookback
statistics are accurate for the first step and progressively worse for the last.
Our \arevin\ keeps RevIN's reversible structure but makes the de-normalization
horizon-aware and gated: a small learnable module modulates the restored
statistics as a function of horizon position, and a gate interpolates between
trusting the lookback statistics and letting the prediction set its own level.
This is philosophically aligned with the de-stationary idea of Non-stationary
Transformers---removed statistics should be re-injected \emph{adaptively} rather
than identically---but is realized as a lightweight, attention-free,
RevIN-compatible module suited to a linear backbone. To our knowledge, a
horizon-adaptive or gated RevIN of this form has not been published, and it is
the component least anticipated by prior linear or frequency-domain models
(Section~\ref{sec:related:positioning}).

\subsection{Frequency-Domain Forecasting}
\label{sec:related:freq}

A large body of work moves part or all of the forecasting computation into a
spectral representation, motivated by the fact that long-term series are
dominated by a few sparse, well-separated periodic components. Autoformer
\citep{wu2021autoformer} introduces the trend/seasonal series-decomposition
block (the moving-average split DLinear later reused) and replaces self-attention
with an Auto-Correlation mechanism that discovers period-based dependencies via
the FFT. FEDformer \citep{zhou2022fedformer} combines series decomposition with
frequency-enhanced attention over a random subset of Fourier (or wavelet) modes,
giving linear complexity and an explicit low-rank spectral filter. FiLM
\citep{zhou2022film} uses Legendre-polynomial projections to compress history and
Fourier projections to denoise it. TimesNet \citep{wu2023timesnet} uses the FFT
to find dominant periods, folds the 1D series into period-indexed 2D tensors, and
applies 2D convolutions. FreTS \citep{yi2023frets} is the most relevant of this
group: it discards attention entirely and learns frequency-domain MLPs over the
real and imaginary parts of the spectrum, across both channel and time
dimensions.

\method\ shares these models' premise---the frequency domain is the right place
to separate predictable structure from noise---but rejects their machinery.
FEDformer, FiLM, and TimesNet are heavy (attention, polynomial bases, or 2D
convolutions), with parameter and FLOP budgets far beyond the small-model regime
a 4\,GB GPU comfortably supports; FreTS is lighter but still uses complex-valued
MLPs across channels and time. \method\ instead uses the spectrum only to
\emph{split} the signal into bands, then forecasts each band with a plain
real-valued linear head in the time domain. The frequency domain is a routing
device, not the space in which the forecast is computed, which keeps the model
linear-scale while still exploiting the spectral separability these works
demonstrate.

\subsection{Transformer, MLP-Mixer, and Multiscale Forecasters}
\label{sec:related:transformer}

At the high-accuracy, high-cost end of the design space, Informer
\citep{zhou2021informer} makes long-sequence forecasting tractable with
ProbSparse attention and a generative decoder; PatchTST \citep{nie2023patchtst}
tokenizes the series into subseries patches and is strictly channel-independent,
and is the strongest Transformer forecaster on the standard benchmarks and our
small-Transformer baseline where memory permits; iTransformer
\citep{liu2024itransformer} inverts the tokenization so each variate becomes a
token and attention models cross-channel dependencies. These motivate the
central question of this paper: how much of their accuracy can a linear-scale
model recover at a fraction of the cost? We adopt PatchTST's channel-independent
treatment but reject patching and attention.

A parallel line removes attention while keeping expressive mixing. TSMixer
\citep{chen2023tsmixer} is an all-MLP architecture alternating time-mixing and
feature-mixing MLPs; TimeMixer \citep{wang2024timemixer} decomposes the series at
multiple sampling scales and mixes the disentangled multiscale components. Both
confirm that decomposition plus simple mixing is a strong, attention-free recipe,
and TimeMixer's multiscale decomposition is a temporal analogue of our spectral
band split. They differ from \method\ in weight class: they stack several mixing
blocks across scales and channels, whereas \method\ uses a single frequency split
with $K$ linear heads and no inter-channel mixing, targeting the smallest viable
model. The efficiency prong is grounded in the Green~AI agenda
\citep{schwartz2020greenai} and the energy-cost analysis of large models
\citep{strubell2019energy}, which argue that efficiency should be reported
alongside accuracy; \method\ reports parameters, FLOPs, train time, and peak
memory as first-class metrics. The Transformer \citep{vaswani2017attention} and
Adam \citep{kingma2015adam} are cited as foundational references.

\subsection{Positioning and Novelty}
\label{sec:related:positioning}

Two recent works are close enough that we position against them explicitly and
honestly. FITS \citep{xu2024fits} is the closest lightweight frequency model: it
applies rFFT to the lookback, a single complex-valued linear layer performing
amplitude scaling and phase shifting, then irFFT, using RevIN and a \emph{fixed
low-pass filter} that \emph{discards} high-frequency components (the source of
its $\sim$10k-parameter footprint), and is channel-independent. The overlap with
\method\ is real---both are RevIN + frequency-domain + lightweight + CI---and
reviewers will rightly raise it. \method\ differs from FITS in three concrete,
testable ways: (i) high-frequency content is \emph{modeled} by a dedicated head
rather than discarded by a low-pass filter, so residual structure that a low-pass
model leaves on the table can be exploited; (ii) the forecast is computed in the
time domain by $K$ real linear heads, not by a single complex map in the
frequency domain---a simpler hypothesis class; and (iii) the normalization is the
horizon-adaptive, gated \arevin, which has no analogue in FITS's plain RevIN.
Accordingly, we do \emph{not} claim to be the first lightweight frequency-domain
linear model---FITS holds that ground---and we include FITS as a run-in-budget
baseline.

FreqMoE \citep{liu2025freqmoe} also decomposes the spectrum into bands and
assigns each to an expert, with a gating network, learnable band boundaries, and
complex-valued residual blocks, using plain instance normalization. It shares the
``each band gets its own head'' idea but is substantially heavier---a
mixture-of-experts with complex-valued multi-block stacks---whereas \method\ uses
$K$ real linear heads recombined by addition, and FreqMoE has no adaptive
normalization contribution. As a 2025 arXiv-only work it is contemporaneous
rather than established prior art; we cite it and differentiate on weight class
and on \arevin. Taken together, the frequency band split with per-band heads is
not unique, so we present it as a specific lightweight design point rather than a
new concept, and we foreground \arevin---the prong with the least prior-art
overlap---as the primary methodological novelty.
